The prespecified loan result replicates (Table~\ref{tab:fresh-primary}). At 64 verified bundles and $q=.99$, augmentation raises mean full-event recovery from 30.2 to 62.5 per block and reduces new damage from 2.08 to 0.00. After averaging the three fixed model/fault runs within each block, recovery improves in all eight blocks. Fitted-seed sensitivity remains: the three seed-specific recovery differences, averaged across blocks, are 5.5, -11.5 and 103.0 events. Thus one fitted-seed mean loses recovery; block uncertainty does not cover new training draws. The paired recovery difference is 32.3 events [95\% interval 29.5, 35.3], and the damage difference is -2.08 [-3.75, -0.75]. H1 meets both its directional and interval criteria. However, history F1 decreases by 0.23 percentage points: its paired difference is -0.23 [-0.30, -0.18]. The confirmed benefit concerns recovery and harm at this operating point, not uniform improvement in every endpoint.

Table~\ref{tab:fresh-primary} also retains the failures. At 256 bundles, augmentation lowers mean history F1 and recovery and increases mean damage in both sources at both targets; not every damage interval excludes zero. On purchasing at 64 bundles it reduces recovery and increases damage at both targets. Its small purchasing F1 gain at 16 bundles and $q=.99$ accompanies fewer full-event recoveries and more damage. Thus generated data can add value beyond repeated corruption, but the evidence supports a narrow operating condition rather than a general verification-budget saving or a universal rule for when to stop generating data.

The never-merge control changes the safety interpretation (Tables~\ref{tab:fresh-safety} and~\ref{tab:fresh-contrasts}). For real-only loans, guarding reduces damage from 38.12 to 2.08, while recovery changes from 30.17 to 30.21. Never merging further reduces damage to 0.29 with the same 30.2 recoveries, but gives up 2.75 true merged pairs and some history F1. On purchasing, guarded real-only/augmented policies retain 9.46/16.08 true pairs per block. Their history F1 exceeds never merge, but their damage is 5.42/16.62, versus 4.71/15.38. The corresponding extra-harm intervals are positive. None of the four H2 comparisons establishes a beneficial merge policy without added harm over never merge.

The strict augmented loan model makes no merges under any of the three policies; all three produce identical outcomes. Its zero observed harm therefore does not demonstrate an incremental guard benefit. Moreover, real-only never merge still recovers 30.2 events versus 62.5 for the augmented model, so the recovery advantage is not explained solely by false merges in the real-only comparator. Never merging eliminates false identity merges by construction, but participation repairs can still damage records. The choice between retaining fragmented histories and accepting residual merge risk remains an application decision; the measured frontier supplies evidence, not a deployment safety guarantee.

All 768 saved repair programs passed independent replay, exact undo, event-content preservation and recomputed history/identity and cluster-expanded damage checks. The source-selection audit verified all older partition hashes and complete coverage of the declared new roots. No new fitting or threshold search was performed. Complete per-block/per-seed outcomes, conditional intervals and measured evaluation costs accompany the manuscript.
